\section{Introduction}\label{sec:introduction}

A stability wall can contract a family of curves and a higher-dimensional
exceptional fibre at the same time. The topology of the smooth moduli
space then depends on how these cycles meet, while its metric degeneration
depends on the relative rates at which their periods vanish. The two
questions concern the same geometric objects. To relate them, a local
description of a singularity must be placed inside a compact moduli
problem, with its full boundary and its actual harmonic metric.

We carry this out for two explicit moduli problems on K3 surfaces. The
first contracts one Lagrangian plane. The second has a moving family of
projective lines and a deeper fibre consisting, in one chamber, of two
Lagrangian planes meeting at a point. A moving line degenerates into one
line in each plane. This specialization gives an integral relation among
global curve classes and an exact identity among their harmonic areas.
The intersection of the planes also changes the integral attachment:
the relevant boundary group is cyclic of order eight, whereas the
boundary of a single plane has order three.

\subsection{The compact family}
For a sheaf on a K3 surface, use the Mukai vector and pairing
\[
 v(F)=(r,c_1,r+c_1^2/2-c_2),\qquad
 (r,c,s)\cdot(r',c',s')=c\cdot c'-rs'-r's.
\]
Choose a K3 surface $X$ with full N\'eron--Severi lattice generated by
$H,D_1,D_2$, where $H$ is ample and the intersection matrix is
\begin{equation}\label{eq:intro-lattice}
 \begin{pmatrix}2&0&0\\0&-6&-4\\0&-4&-8\end{pmatrix}.
\end{equation}
Such a surface exists by the construction in
Section~\ref{sec:compact-interaction}. Put
\[
 L_0=\OO_X,\quad L_1=\OO_X(D_1),\quad L_2=\OO_X(D_2),
 \qquad a_i=v(L_i),\qquad v=a_0+a_1+a_2.
\]
Thus $v=(3,D_1+D_2,-4)$ and $v^2=2$. Let $\omega_\eta$ be the
Ricci-flat representative of $H+sD_1+tD_2$, for
$\eta=(s,t)$ near zero. The normalized stability parameters are
\begin{equation}\label{eq:intro-levels}
 \zeta(\eta)=\frac{2\pi}{3}
 \begin{pmatrix}-10&-12\\8&0\\2&12\end{pmatrix}
 \binom{s}{t},\qquad \zeta_0+\zeta_1+\zeta_2=0.
\end{equation}
The three lines $\zeta_i=0$ divide the parameter plane into six
chambers. By a generic parameter below we mean a point outside these
lines in a sufficiently small neighborhood of the origin.

On a smooth moduli space we use the real part of the unnormalized
Dolbeault harmonic $L^2$ pairing, at fixed holomorphic determinant.
Write $g_\eta$ for this metric and $\Omega_\eta$ for its K\"ahler
form. All distances between exceptional points are measured in the
entire moduli space. Integral cohomology is understood unless another
coefficient group is displayed.

\begin{theorem}\label{thm:main}
For the family above, the following assertions hold.
\begin{enumerate}
\item Each generic chamber $\sigma$ has a smooth compact projective
moduli fourfold $M_\sigma$ of slope-stable bundles, of deformation
type $X^{[2]}$. There are projective symplectic contractions
$\pi_\sigma:M_\sigma\to Z$ to a common normal projective fourfold.
The regular locus of $Z$ is the moduli of $H$-stable bundles. Its
reduced singular locus is a normal K3 surface $T$ with one $A_1$ point $o$.
The fibres over $T\setminus\{o\}$ are projective lines. Over $o$
the fibre is two planes meeting transversely at a point in the
chambers with signs $(+,-,+)$ and $(-,+,-)$; in the other four
chambers it is $\PP^2\cup\FF_1$, meeting along a line and the
negative section. Every central or nearby ideal HYM object has zero
lost charge, and every sheaf in these moduli problems is locally free.

\item Fix $\zeta_0=a>0$, $\zeta_2=b>0$ and
$\zeta_1=-a-b$. Write $P_0,P_2$ for the two planes, $c_0,c_2$
for their line classes, and $f$ for a moving fibre class. These
classes satisfy
\begin{equation}\label{eq:intro-specialization}
 f=c_0+c_2,
 \qquad
 \left(\int_{M_\sigma}[P_i][P_j]\right)_{i,j=0,2}
   =\begin{pmatrix}3&1\\1&3\end{pmatrix}.
\end{equation}
The curve and surface pairs each span a primitive rank-two integral
subgroup. A proper flat analytic family realizes the curve relation
with both multiplicities equal to one. For
$V=M_\sigma\setminus(P_0\cup P_2)$, let $L$ be the boundary of
a regular neighborhood of the planes. Then $H^4(L)=\Z/8$, and
restriction to $L$ is an explicit surjection
$\kappa:H^4(V)\to\Z/8$ such that restriction identifies
\begin{equation}\label{eq:intro-attachment}
 H^4(M_\sigma)=
 \{(u,n_0,n_2)\in H^4(V)\oplus\Z^2:
       n_0-3n_2\equiv\kappa(u)\pmod 8\}.
\end{equation}
Proposition~\ref{prop:calculated-attachments} evaluates the map on
Chern classes of family indices associated with a universal bundle.
The complement cohomology ring, together with these distinguished
restricted classes, is retained global data.

\item The same cycles have exact harmonic areas and volumes
\begin{equation}\label{eq:intro-volumes}
 \begin{gathered}
 \operatorname{Area}(c_0)=\pi a,\qquad
 \operatorname{Area}(c_2)=\pi b,\qquad
 \operatorname{Area}(f)=\pi(a+b),\\
 \Vol(P_0)=\frac{\pi^2a^2}{2},\qquad
 \Vol(P_2)=\frac{\pi^2b^2}{2}.
 \end{gathered}
\end{equation}
Their ambient diameters are comparable to $\sqrt a$ and $\sqrt b$,
with constants independent of $a/b$. As $a/b\to0$ and
$a+b\to0$, rescaling distances by $a^{-1/2}$ at $P_0$ gives the
complete unit-level Calabi quotient on $T^*\PP^2$, and
\[
 \frac{\diam_{M_\sigma}P_0}{\sqrt a}\longrightarrow\frac\pi2.
\]
The reversed hierarchy has the corresponding conclusion for $P_2$.
There are also complete successive pointed limits at any fixed finite
number of scales, under the convergence hypotheses on normalized
levels and normal basepoints in Theorem~\ref{thm:successive-limits}.
The harmonic metric on $Z_{\rm reg}$ has compact completion
$(Z,d_0)$ with its projective topology, and
\[
 (M_{\sigma(\eta)},g_\eta)\longrightarrow(Z,d_0)
 \quad\text{in compact Gromov--Hausdorff distance}
\]
along every generic sequence $\eta\to0$.
\end{enumerate}
\end{theorem}

The formula \eqref{eq:intro-attachment} records an integral additive
attachment. Its order-eight congruence depends on the intersection
of the two planes. It does not follow from the ranks of their
cohomology groups. The complement $V$ still contains the exceptional
geometry over the moving stratum; it differs from the common stable
open set. Section~\ref{sec:projective-family} describes that geometry
and computes its ring over the open stratum, retaining the nonzero
algebraic Brauer class of its projective-line bundle.

\subsection{A single plane and the proof mechanism}
Section~\ref{sec:isolated-plane} begins with the lattice
$\NS(X)=\Z H\oplus\Z D$, $H^2=2$, $H\cdot D=0$, $D^2=-22$,
and vector $(3,D,-4)$. The only strictly polystable central bundle
is $\OO_X\oplus B$, with $B$ a rigid rank-two bundle. Each adjacent
compact fourfold contracts a plane. Its primitive surface class has
square three, and the boundary of a normal neighborhood has fourth
cohomology $\Z/3$. The corresponding generator formula for the
attachment is the one-component version of
\eqref{eq:intro-attachment}. This example also fixes the metric
normalization: if $\kappa_t=44\pi|t|/3$, an exceptional line has
area $\pi\kappa_t$, and the plane has ambient diameter
$(\pi/2)\sqrt{\kappa_t}(1+O(|t|))$.

Three arguments connect these calculations to compact geometry.
First, an exhaustive rank calculation uses the lower bound $-2$ for
the Mukai square of a stable K3 sheaf to exclude every possible
lost-charge vector. Surface gauge compactness then applies without
concentration. Normalized inclusions and quotients in the central
Jordan--H\"older filtrations show that entire exceptional fibres
enter the local analytic charts. This supplies the global coverage
needed before a quotient model can identify a contraction.

Second, universal extension families determine restrictions of
integral classes to the exceptional cycles. These restrictions prove
primitivity and evaluate the boundary homomorphism. The family
$xy=t$ at the intersection of the two planes identifies the actual
specialization multiplicities. The two ingredients serve different
purposes: lattice pairings identify classes, while a proper flat
family proves their geometric specialization.

Third, the harmonic comparison is relative in every quotient
direction. Its error remains small even when one period is much
smaller than the other. A lower bound for the length required to
leave a chart then compares the two ambient distance infima. This
explains why the estimates apply to distances in the whole compact
moduli space, and why the rescaled limits are complete. Global
completion additionally uses the identification of the limiting
harmonic class with the ample class of the projective contraction.

\subsection{Relation to earlier work}
The existence and deformation type of the compact moduli spaces are
provided by Yoshioka's theorems \cite{yoshioka99,yoshioka01}. We use
the projective slope moduli of Greb--Toma \cite{grebtoma} and the
surface gauge comparison of Greb--Sibley--Toma--Wentworth
\cite{gstw}. Bayer--Macr\`i's theory \cite{bayermacri} supplies
the broader birational setting for stability walls. Local quiver
descriptions of K3 sheaf moduli are developed by Arbarello--Sacc\`a
\cite{arbarellosacca}. The two compact fibre configurations in
Theorem~\ref{thm:main} already occur in Bielawski--Dancer's
hypertoric examples \cite[Example~6.10]{bielawskidancer}; here they
are realized as the complete fibres of the stated compact
slope-moduli contractions.

Formality is also an established input to this subject.
Bandiera--Manetti--Meazzini \cite[Theorem~5.1]{bmm} prove DG-Lie
formality for coherent sheaves with linearly reductive automorphism
group on a smooth projective surface with trivial canonical bundle.
Their theorem applies to our common-slope direct sums. The related
projectively hyperholomorphic formality theorem is
\cite[Theorem~4.3]{meazzinionorati}. Section~\ref{sec:harmonic-models}
gives the additional coordinate information used for the metric
comparison: a convergent equivariant symplectic change whose first
correction is cubic. Relative finite balancing and the actual
harmonic reduction are the previously proved results
\cite{wuebben14,wuebben15}, restated below with the required
normalizations.

The plane-class identities and Hilbert-square cohomology are
classical \cite{hassetttschinkel,totaro,gottsche}; the calculation
here identifies their integral representatives and attachments in
the compact bundle-moduli problem. Markman's index generators
\cite{markman} make the attachment maps explicitly evaluable.
For the global metric completion we use Song \cite{song} and
Tosatti \cite{tosatti}, after checking their geometric and
metric hypotheses. The conclusions about cycle diameters and
successive scales have separate local proofs.

The resulting theorems concern exceptional cycles of compact
bundle-moduli contractions without bubbling. The complement and its
distinguished integral classes supply the global information retained
in the cohomological calculation.

